\chapter{Conclusion and Future Work}
\label{sec:conclusion}

This chapter should close the thesis by restating the problem, summarizing the main findings, and identifying the next research steps. It should remain tightly aligned with the actual evidence presented in Chapter~\ref{sec:results}.

\section{Conclusion}

\begin{itemize}
\item Restate the motivation for tool-orchestrated agentic reverse engineering.
\item Summarize the system contribution.
\item Summarize the main empirical takeaway without overstating certainty.
\item Distinguish clearly between what was learned from preliminary development runs and what was established by the final held-out evaluation.
\end{itemize}

\begin{itemize}
\item \TODO{Write a concise closing argument that answers the thesis research questions directly.}
\item \TODO{State the strongest defensible claim supported by the final held-out experiments, not just by the earlier development or pilot runs.}
\item \TODO{Add one sentence that explicitly limits the interpretation of the preliminary broad-coverage and decoder-depth runs to hypothesis-shaping and methodology refinement where appropriate.}
\item \TODO{State clearly that the held-out malware run was designed to keep runtime family identity blind and to prioritize evidence recovery over family naming.}
\end{itemize}

\section{Future Work}

\begin{itemize}
\item Broader and more realistic malware corpora.
\item Stronger human evaluation and analyst-in-the-loop studies.
\item Richer dynamic-analysis orchestration and sandbox integration.
\item More principled automatic artifact selection, including packed versus unpacked analysis targets.
\item Improved cost-aware planning, recovery, and significance analysis.
\end{itemize}

\begin{itemize}
\item \TODO{Prioritize the future-work items that best follow from the actual limitations chapter.}
\item \TODO{Call out which extensions are engineering scale-ups versus genuinely new research questions.}
\end{itemize}
